\begin{table*}[t]
\centering
\caption{Posterior-risk posture in matched TRACE-BiOpt route probes.}
\label{tab:trace-biopt-posterior-risk}
\footnotesize
\begin{tabularx}{\textwidth}{llcccc>{\raggedright\arraybackslash}X}
\toprule
Case & Comparator & $\Delta$ trace & $\Delta$ MAE & Best-route posterior share & Best route & Reading \\
\midrule
Seattle 20\% & Zero-weight probe & 20.4 & 0.036 & 63.1\% & Seattle diag. certified & Lower trace aligns with lower held-out MAE in this matched probe. \\
Seattle 30\% & Zero-weight probe & 9.4 & 0.032 & 63.8\% & Seattle diag. certified & Lower trace aligns with lower held-out MAE in this matched probe. \\
PeMS7\_228 10\% & Stage15 certified & 35.6 & 0.129 & 18.8\% & Calibrated low-cert & Lower trace aligns with lower held-out MAE in this matched probe. \\
PeMS7\_228 10\% & Zero-weight probe & 28.1 & 0.169 & 18.8\% & Calibrated low-cert & Lower trace aligns with lower held-out MAE in this matched probe. \\
\bottomrule
\end{tabularx}
\begin{flushleft}\footnotesize $\Delta$ trace and $\Delta$ MAE are comparator minus the reference TRACE-BiOpt route in the same case, so positive values mean the reference route has both lower posterior trace and lower held-out MAE. The Seattle rows are bounded same-split diagnostic slices, while the PeMS7\_228 row uses the promoted calibrated current-best route. This is a bounded theorem-to-evidence bridge for Theorem~\ref{thm:trace-risk}: it uses only matched route probes inside the same dataset-budget slice and does not claim that posterior trace must rank every audited baseline by MAE under non-Gaussian traffic data.\end{flushleft}
\end{table*}
